% Synthetic PRD-style manuscript for segmenter tests (paper-copilot).
% It mixes the constructs a real revtex4-2 draft uses. Periods in comments. Must be ignored.
\documentclass[aps,prd,twocolumn,superscriptaddress,nofootinbib]{revtex4-2}
\usepackage{amsmath,graphicx,hyperref}
\usepackage[capitalise]{cleveref}
\newcommand{\be}{\begin{equation}}
\newcommand{\ee}{\end{equation}}
\newcommand{\etal}{et al.}
\newcommand{\pT}{p_{\mathrm{T}}}
\newcommand{\ttbar}{\ensuremath{t\bar{t}}}
\newcommand{\GNN}{GNN}
\def\HL{HL-LHC}

\begin{document}

\preprint{UCI-TR-2026-01}
\title{Tracking displaced particles with a graph neural network trained on generic events}
\author{A. Author}
\email{a.author@uci.edu}
\thanks{Corresponding author. Second sentence of a thanks.}
\affiliation{Department of Physics and Astronomy, University of California, Irvine, CA 92697, USA}
\author{B. Author}
\affiliation{CERN, 1211 Geneva 23, Switzerland}
\collaboration{The Canto Collaboration}
\date{\today}

\begin{abstract}
Long-lived particles (LLPs) decay far from the collision point. Their tracks are hard to find, e.g.\ for $d_0 > 10$~mm.
We train a \GNN{} tracker on \ttbar{} events only and test it on displaced signals.
It finds 96.9\% of the targets, i.e., 2.3 points more than the baseline. Is this enough?
\end{abstract}

\maketitle

\section{Introduction}
\label{sec:intro}
Searches for LLPs at the \HL{} need efficient tracking~\cite{Ju2022Zenodo,Lee2019}.
Graph neural networks were shown to work for prompt tracks by Ju \etal~\cite{Ju2021}, and Wang et al.\ extended them to large radii~\cite{Wang2022LRT}.
这里写一下为什么只用ttbar训练很重要
However, every study so far trained on the signal it evaluates (cf.\ Refs.~\cite{Lee2019,Wang2022LRT}).
In this paper we ask a simple question: does a tracker trained without the signal still find it?
We show that it does.\footnote{A preliminary version appeared in a workshop note. It used fewer events.} The answer holds for all species (see Sec.~\ref{sec:results}).

The rest of this paper is organized as follows.
Section~\ref{sec:method} defines the targets; \cref{sec:results} gives the results; App.~\ref{app:sim} describes the simulation.
$N$ events are used in total.
1000 events are kept for validation.
\Cref{fig:eff} summarizes the main result.

\section{Method}
\label{sec:method}
A target particle has $\pT > 1$~GeV and $|\eta| < 4$ (approx.\ the tracker acceptance).
The efficiency is defined as
\be
  \epsilon = \frac{N_{\text{matched}}}{N_{\text{targets}}}.
\ee
The fake rate is
\begin{align*}
  f &= \frac{N_{\text{fake}}}{N_{\text{tracks}}},
\end{align*}
where a fake track matches no particle.
Both rates use the same population, cf.~Eq.~(\ref{eq:eff}) and Eqs.~\eqref{eq:a}--\eqref{eq:b}.
The total uncertainty combines two terms,
\[
  \sigma^2 = \sigma_{\text{boot}}^2 + \sigma_{\text{seed}}^2 .
\]
The bootstrap term dominates. % a comment. With a period.
We quote 50\% and 95\% intervals; the price is \$0 per event (vs.\ \$5 for the full simulation).

\paragraph{Selection.} Events pass a trigger with $\pT > 25$~GeV. We then apply
\begin{itemize}
  \item a vertex cut, and
  \item a track-quality cut. It removes 3\% of the tracks.
\end{itemize}
The cuts are listed in \cref{tab:cuts}.

\begin{description}
  \item[Efficiency:] the fraction of targets with a matched track.
  \item[Fake rate:] the fraction of tracks without a matched target.
\end{description}

\subsection{Training}
The model is trained for 50 epochs with a learning rate of $10^{-3}$.
We use three seeds. The spread among them is small.
``Generic'' here means \emph{no signal at all}. The \textbf{only} input is the hit list.

\section{Results}
\label{sec:results}
On isolated particles the tracker finds $96.94\%$ of the targets, compared with $94.62\%$ for the baseline~\cite{Ju2021}.
This is a gain of $2.32 \pm 0.28$ points (\cref{fig:eff}).
The gain is largest at large displacement, w.r.t.\ the prompt region.
In ttbar the results are similar; see Figs.~\ref{fig:eff} and~\ref{fig:fake}.

\begin{figure}[t]
  \centering
  \includegraphics[width=\linewidth]{fig/eff.pdf}
  \caption[Efficiency versus displacement]{Efficiency versus transverse displacement. Solid: the \ttbar-trained tracker; dashed: the baseline. The ratio is shown in the lower panel.}
  \label{fig:eff}
\end{figure}

\begin{figure*}
  \subfloat[Efficiency.]{\includegraphics[width=0.45\linewidth]{a.pdf}\label{fig:sub-a}}
  \subfloat[Fake rate.]{\includegraphics[width=0.45\linewidth]{b.pdf}}
  \caption{Comparison of the two trackers. Both use the same track builder.}
  \label{fig:fake}
\end{figure*}

\begin{table}[b]
  \caption{Selection cuts.\label{tab:cuts}}
  \begin{tabular}{lc}
    \hline
    Cut & Value. With a period. \\
    $\pT$ & $> 1$~GeV~\cite{Lee2019} \\
    \hline
  \end{tabular}
\end{table}

\iffalse
This paragraph was commented out. It must not appear.
\fi

\begin{verbatim}
Verbatim text. It must not appear.
\end{verbatim}

\section{Summary and outlook}
\label{sec:concl}
We have shown that generic training is enough.
Future work will add pileup.

\begin{acknowledgments}
We thank the Canto team for discussions. This work was supported by the U.S. Department of Energy.
\end{acknowledgments}

\appendix
\section{Detector simulation}
\label{app:sim}
The detector is simulated with Geant4~\cite{Geant4}.
The magnetic field is 2~T.

\bibliography{refs}

\end{document}
Text after the end of the document. It must not appear.
